\documentclass[a4paper,twoside]{article}

\usepackage{morewrites}

\usepackage[svgnames,dvipsnames,x11names]{xcolor}
\usepackage{graphicx}
\usepackage[english]{babel}
\usepackage[colorlinks=true, linkcolor=NavyBlue, urlcolor=NavyBlue, citecolor=NavyBlue]{hyperref}
\usepackage[a4paper]{geometry}
\usepackage{ts-fonts}
\usepackage{tcolorbox}
\usepackage{fvextra}
\usepackage{amsmath}
\usepackage{float}
\usepackage{seqsplit}
\usepackage{ts-code}

\usepackage{lastpage}
\usepackage{refcount}
\makeatletter
\def\ps@plain{
\let\@oddhead\@empty
\let\@evenhead\@empty
\def\@oddfoot{
\hfil
\ifnum\getpagerefnumber{LastPage}>1\relax
\thepage
\fi
\hfil}
\let\@evenfoot\@oddfoot
}
\makeatother

\title{Glossary}
\author{}\date{}
\begin{document}
\maketitle
The MkDocs \href{https://realtimeprojects.github.io/mkdocs-ezglossary}{ezglossary} plugin exists, but the syntax is unintuitive (for example, the semicolon separator) and does not handle spaced or formatted entries well. Ideally, readers should be able to access the glossary as a popup, tooltip, or link.

\LaTeX{} glossaries usually rely on the \tscodeinline{glossaries} package and definitions declared in the preamble via \tscodeinline{\textbackslash{}newglossaryentry}:
\begin{tscode}[lang=latex, engine=pygments]
\begin{Verbatim}[commandchars=\\\{\},breaklines, breakanywhere, commandchars=\\\{\}]
\PY{k}{\PYZbs{}newglossaryentry}\PY{n+nb}{\PYZob{}}key\PY{n+nb}{\PYZcb{}}\PY{n+nb}{\PYZob{}}
  name=\PY{n+nb}{\PYZob{}}singular form\PY{n+nb}{\PYZcb{}},
  plural=\PY{n+nb}{\PYZob{}}plural form\PY{n+nb}{\PYZcb{}},
  description=\PY{n+nb}{\PYZob{}}descriptive text\PY{n+nb}{\PYZcb{}},
  first=\PY{n+nb}{\PYZob{}}special form for first use\PY{n+nb}{\PYZcb{}},
  text=\PY{n+nb}{\PYZob{}}normal form (if you want to force it)\PY{n+nb}{\PYZcb{}},
\PY{n+nb}{\PYZcb{}}
\end{Verbatim}
\end{tscode}
\tscodeinline{\textbackslash{}gls} uses the standard form, \tscodeinline{\textbackslash{}Gls} capitalizes the first letter, \tscodeinline{\textbackslash{}glspl} gives the plural, and \tscodeinline{\textbackslash{}Glspl} capitalizes the plural. The first use of the term falls back to the \tscodeinline{first} form when present, otherwise \tscodeinline{name}. Markdown alone cannot express all of that. One option is to define glossary entries in front matter:
\begin{tscode}[lang=yml, engine=pygments]
\begin{Verbatim}[commandchars=\\\{\},breaklines, breakanywhere, commandchars=\\\{\}]
glossary:
  html:
    name: HTML
    plural: HTMLs
    description: \PYZgt{}
      HyperText Markup Language, the standard markup language for
      documents designed to be displayed in a web browser.
    first: HyperText Markup Language (HTML)
\end{Verbatim}
\end{tscode}
Place an anchor where the term is defined:
\begin{tscode}[lang=markdown, engine=pygments]
\begin{Verbatim}[commandchars=\\\{\},breaklines, breakanywhere, commandchars=\\\{\}]
The \PY{g+gs}{**HTML**}\PYZob{}\PYZsh{}gls\PYZhy{}html\PYZcb{} is the standard markup language for web pages...
\end{Verbatim}
\end{tscode}
\end{document}
